\section{Related Work}

Semantic caching originated in relational database systems, where Dar et
al.~\cite{dar1996semantic} introduced cached query results and semantic
containment checks to reduce query-processing cost. For LLM serving, the
dominant architecture translates this concept to vector spaces. GPTCache
\cite{bang2023gptcache} embeds incoming prompts, retrieves the nearest
cached embedding by cosine similarity, and serves the stored response if
the similarity clears a single configured threshold. MeanCache
\cite{gill2025meancache} targets a federated, user-centric deployment
setting but retains the identical univariate fixed-threshold gate.

vCache \cite{schroeder2026vcache}, the strongest learned baseline,
observes that a global threshold is miscalibrated across heterogeneous
embedding neighborhoods. It replaces the global cutoff with per-embedding
thresholds calibrated online from local similarity statistics. This
innovation improves calibration, but it does not alter \emph{which
evidence} drives the decision: the reuse verdict still depends on top-1
similarity alone. In all three systems GPTCache, MeanCache, and vCache,
the top-$k$ nearest-neighbor search produces several retrieval-geometry
signals: the runner-up similarity, the margin between the top-1 and
runner-up matches, and the local embedding density around the query. All
three signals are computed and then discarded.

This uniformity reveals a structural gap. Every existing semantic cache
for LLMs treats gating as thresholding on a single feature. The decision
boundary is a scalar cutoff, whether it is global, per-embedding, or
hand-tuned. No prior system casts cache gating as a binary classification
problem amenable to supervised learning. No system exploits the richer
retrieval-geometry features that the cache already produces. The notion
of a learned multi-feature gate that predicts reuse correctness from the
full retrieval signature, while leaving the cache index and eviction
policy unchanged, has not been proposed.

SERVE fills this gap with a light-weight gradient-boosted tree ensemble
that operates over four zero-cost features: top-1 similarity, margin,
runner-up similarity, and local density. The gate intercepts the cache's
retrieval result, scores the feature vector, and outputs a binary decision
node by node. The features are byproducts of the existing nearest-neighbor
lookup, so the gate adds no further embedding calls or index operations to
the retrieval pipeline. The ensemble architecture provides sub-microsecond
inference latency and explicit, inspectable decision rules. The gate is a
standalone component compatible with any semantic cache backend, and
this architectural decoupling, together with the supervised classification
formulation, distinguishes the approach from prior caching systems.

A second, orthogonal gap concerns query recurrence. No prior work on
semantic caching analyzes how gating effectiveness varies with the
frequency at which queries reappear in the workload. SERVE's relative
hit-rate gain over a fixed-threshold gate at a 1\% false-serve budget
grows monotonically with recurrence, from $+21.3\%$ at $r=0.2$ to
$+81.0\%$ at $r=0.9$, where $r$ is the proportion of test queries that
draw from the cache.
% E6
Recurrence-aware evaluation exposes a fresh dimension for designing and
benchmarking cache gating mechanisms, and it explains why the largest
absolute gains concentrate in the high-recurrence, strict-budget regime
where even modest per-query improvements compound over many repeated
serves.
